%%
%% Boost.Decimal: outline of a research paper for ACM Transactions on
%% Mathematical Software (TOMS).
%%
%% This is a complete first draft: every section is written in prose. The
%% remaining camera-ready placeholders are the rights/DOI/volume block, the CCS
%% concept codes, the received dates, and final verification of references.bib.
%%
%% It uses the unmodified ACM template in ../template (acmart.cls and
%% ACM-Reference-Format.bst). Build with the provided Makefile, which points
%% LaTeX at that folder without changing it. See README.md.
%%
%% TOMS is published in the "acmsmall" journal format.
%%
\documentclass[acmsmall]{acmart}

\AtBeginDocument{%
  \providecommand\BibTeX{{%
    Bib\TeX}}}

%% Code listings: compact, monochrome C++ style. Standard listings package.
\usepackage{listings}
\lstset{
  language=C++,
  basicstyle=\footnotesize\ttfamily,
  columns=fullflexible,
  keepspaces=true,
  showstringspaces=false,
  breaklines=true,
  frame=single,
  framesep=4pt,
  xleftmargin=4pt,
  belowcaptionskip=4pt,
  aboveskip=6pt,
  belowskip=4pt
}

%% Rights: SAMPLE values. Replace with the values from the ACM rights form
%% once the paper is accepted.
\setcopyright{acmlicensed}
\copyrightyear{2026}
\acmYear{2026}
\acmDOI{XXXXXXX.XXXXXXX}

%% Journal: TOMS.
\acmJournal{TOMS}
\acmVolume{0}
\acmNumber{0}
\acmArticle{0}
\acmMonth{0}

%% Title. Keep a short title for the running header.
\title[Boost.Decimal: Portable IEEE 754 Decimal Floating-Point]{Boost.Decimal:
  A Portable, Header-Only, constexpr and GPU-Capable Implementation of
  IEEE 754 Decimal Floating-Point Arithmetic}

%% Authors. ORCID for Matt Borland is taken from the JOSS submission.
\author{Matt Borland}
\orcid{0009-0005-8183-4254}
\affiliation{%
  \institution{The C++ Alliance}
  \country{USA}}
\email{matt@mattborland.com}

\author{Christopher Kormanyos}
\affiliation{%
  \institution{Independent Researcher}
  \country{Germany}}
\email{e\_float@yahoo.com}

\renewcommand{\shortauthors}{Borland and Kormanyos}

%%
%% Abstract. Draft below; tighten once results are final.
\begin{abstract}
  Decimal floating-point arithmetic, standardized in IEEE 754, represents
  values such as 0.1 exactly and is required wherever results must match
  base-ten expectations, most notably in finance and commerce. The two
  established implementations, the Intel Decimal Floating-Point Math Library
  and IBM decNumber, are C libraries: they are not header-only, expose no
  C++ compile-time interface, do not run on accelerators, and do not cover
  the full range of CPU architectures used today. We present Boost.Decimal,
  a header-only C++14 implementation of the IEEE 754 decimal32, decimal64,
  and decimal128 types (and three non-conforming variants tuned for runtime
  speed) that requires no external dependency, evaluates at compile time
  through constexpr, executes on CPUs and on NVIDIA GPUs from a single source,
  and is validated for cross-platform bit-for-bit consistency across x86,
  ARM, s390x, PPC64LE, and bare-metal targets. We describe the design, the
  arithmetic algorithms, the conformance and reproducibility testing, and a
  performance study showing that Boost.Decimal matches or exceeds the Intel
  reference library across the basic operations on several platforms, for example
  computing decimal64 additions several times faster than the Intel reference on
  x86-64. Because the types satisfy the concept requirements of generic C++ numerical
  code, existing libraries such as Boost.Math operate on decimal data without
  modification, making Boost.Decimal a basis for reproducible decimal analysis. The
  library is available at \url{https://github.com/boostorg/decimal}.
\end{abstract}

%%
%% CCS concepts. PLACEHOLDER: regenerate the exact codes with the ACM CCS
%% tool at https://dl.acm.org/ccs and paste the generated CCSXML here.
\begin{CCSXML}
<ccs2012>
 <concept>
  <concept_id>10002950.10003705</concept_id>
  <concept_desc>Mathematics of computing~Mathematical software</concept_desc>
  <concept_significance>500</concept_significance>
 </concept>
 <concept>
  <concept_id>10011007.10011006.10011073</concept_id>
  <concept_desc>Software and its engineering~Software libraries and repositories</concept_desc>
  <concept_significance>300</concept_significance>
 </concept>
 <concept>
  <concept_id>10002950.10003714</concept_id>
  <concept_desc>Mathematics of computing~Numerical analysis</concept_desc>
  <concept_significance>300</concept_significance>
 </concept>
 <concept>
  <concept_id>10010583.10010662.10010674</concept_id>
  <concept_desc>Computing methodologies~Parallel computing methodologies</concept_desc>
  <concept_significance>100</concept_significance>
 </concept>
</ccs2012>
\end{CCSXML}

\ccsdesc[500]{Mathematics of computing~Mathematical software}
\ccsdesc[300]{Software and its engineering~Software libraries and repositories}
\ccsdesc[300]{Mathematics of computing~Numerical analysis}
\ccsdesc[100]{Computing methodologies~Parallel computing methodologies}

\keywords{decimal floating-point, IEEE 754, decimal32, decimal64, decimal128,
  binary integer decimal, BID, C++, constexpr, header-only library, GPU, CUDA,
  portable numerical software, reproducibility, Boost}

%% Received dates: filled in at acceptance.
%\received{XX Month 2026}
%\received[revised]{XX Month 2026}
%\received[accepted]{XX Month 2026}

\begin{document}

\maketitle

%% ===================================================================
\section{Introduction}
\label{sec:intro}

The binary floating-point types that dominate scientific computing store the
significand in base two. As a consequence they cannot represent most finite
base-ten fractions exactly: the value 0.1 has no finite binary expansion, and the
familiar evaluation of \texttt{0.1 + 0.2} in IEEE 754 binary64 yields
0.30000000000000004 rather than 0.3.\footnote{\url{https://0.30000000000000004.com}}
The error is small, but it is present before any arithmetic is performed, because
it enters when a base-ten literal or an input string such as a price is rounded to
the nearest binary value. In domains where a computed result must equal the value a
person obtains by hand, in particular finance, accounting, tax and currency
calculations, and commerce, this discrepancy is unacceptable, and it is the reason
the 2008 revision of IEEE 754 standardized decimal floating-point arithmetic
\cite{IEEE754-2008, Cowlishaw2003}.

\begin{lstlisting}[caption={Binary \texttt{double} rounds the literals 0.1 and 0.2; the
  decimal type represents them exactly.},label={lst:point}]
using namespace boost::decimal::literals;
std::cout << 0.1 + 0.2 << '\n';        // double:      0.30000000000000004
std::cout << 0.1_DD + 0.2_DD << '\n';  // decimal64_t: 0.3
\end{lstlisting}

Decimal floating-point types store the significand in base ten, so a value such as
0.1 or 0.3 is represented exactly, and they offer a far wider range than the
fixed-point representations that are sometimes used to avoid binary rounding.
Despite the standard being more than fifteen years old, decimal arithmetic remains
awkward to use from modern C++. The two established implementations, the Intel
Decimal Floating-Point Math Library \cite{cornea2009, intel-dfp} and IBM's decNumber
\cite{ibm-decnumber}, are C libraries. They must be compiled and linked rather than
simply included, they expose no compile-time (\texttt{constexpr}) interface, they do
not execute on graphics processors, and neither covers the full range of CPU
architectures in use today. The C++ standards proposals that would have added
decimal types \cite{tr24733, wg21-n3407, wg21-n3871} planned to wrap these C
libraries and left \texttt{constexpr} support optional, and none has been adopted.

This article presents Boost.Decimal, a from-scratch implementation of IEEE 754
decimal floating-point arithmetic that closes this gap. It provides the three
standard types \texttt{decimal32\_t}, \texttt{decimal64\_t}, and
\texttt{decimal128\_t} (7, 16, and 34 significant decimal digits), together with
three companion types that relax the storage layout for greater runtime speed. The
library is part of the Boost C++ Libraries collection, distributed since release
1.91.\footnote{\url{https://github.com/boostorg/decimal}} Its contributions are:

\begin{enumerate}
\item A header-only, dependency-free implementation that requires only C++14 and
  builds with toolchains as old as \texttt{clang++ 6}, removing the integration
  friction of the existing C libraries (Section~\ref{sec:design}).
\item Full \texttt{constexpr} evaluation: every operation, including parsing,
  formatting, and the elementary functions, can be evaluated at compile time
  (Sections~\ref{sec:design} and \ref{sec:impl}).
\item Execution on NVIDIA GPUs through CUDA, from the same source as the CPU
  path. To our knowledge this is the first decimal floating-point library with
  GPU support (Section~\ref{sec:gpu}).
\item Verified bit-for-bit identical results across a wide matrix of operating
  systems and architectures, including big-endian s390x and bare-metal ARM
  Cortex-M (Sections~\ref{sec:portability} and \ref{sec:verification}).
\item A performance study showing that Boost.Decimal matches or exceeds the Intel
  reference library across the basic operations, and is faster than the GCC
  \texttt{\_Decimal} built-in types in many of them (Section~\ref{sec:performance}).
\item Composition with the existing C++ numerical ecosystem: because the types model
  the concepts that generic numerical code expects, libraries such as Boost.Math operate
  on decimal data unchanged, putting a broad range of analysis within reach on exact
  decimal values (Section~\ref{sec:applications}).
\end{enumerate}

We are deliberate about what is and is not new. The decimal formats, the binary
integer significand (BID) encoding, and conformance to IEEE 754 are established
prior art, shared with the Intel and IBM libraries. The contribution of this work
is the combination that no prior implementation offers: a single portable,
header-only, \texttt{constexpr} source that runs conformantly on CPUs, GPUs, and
bare-metal targets while matching the performance of the reference software.

%% ===================================================================
\section{Background and Related Work}
\label{sec:background}

\subsection{IEEE 754 decimal floating-point}
The 2008 revision of IEEE 754 \cite{IEEE754-2008}, retained in the 2019 revision
\cite{IEEE754-2019}, defines three interchange formats for decimal floating-point:
decimal32, decimal64, and decimal128, providing 7, 16, and 34 significant decimal
digits respectively. A finite value is a signed integer significand scaled by a
power of ten. The standard specifies the arithmetic operations, the rounding modes
(with round to nearest, ties to even as the default), the special values (signed
zeros, infinities, and quiet and signaling NaNs), and the behavior of conversions,
following the General Decimal Arithmetic specification of Cowlishaw
\cite{Cowlishaw2003}.

A property unique to decimal is the cohort. Because the significand and exponent are
stored separately, one numerical value has several representations: for example 1e5,
10e4, and 100e3 all denote one hundred thousand. Binary formats have no cohorts.
Preserving or normalizing cohorts is a design concern that does not arise for binary
types, and we return to it when discussing input and output (Section~\ref{sec:io}).

\subsection{Encodings: BID and DPD}
IEEE 754 permits two encodings of the significand field. The Binary Integer Decimal
(BID) encoding stores the significand as a binary integer and is intended for
software implementations; the Densely Packed Decimal (DPD) encoding packs three
decimal digits into ten bits and is intended for hardware. Boost.Decimal stores
values in BID internally for the conforming types, and provides conversions to and
from both BID and DPD bit strings, so values can be exchanged with hardware decimal
units and with other implementations.

\subsection{Existing implementations}
The reference software implementation of decimal arithmetic is the Intel Decimal
Floating-Point Math Library, which Cornea et al. introduced as the first software
implementation of IEEE 754 decimal using the BID encoding \cite{cornea2009}; it is
the de facto industry standard \cite{intel-dfp}. IBM's decNumber is a portable ANSI C
reference implementation of General Decimal Arithmetic and is used inside GCC and the
International Components for Unicode \cite{ibm-decnumber}. GCC also ships libdecnumber
and a copy of the Intel BID library to support its \texttt{\_Decimal32},
\texttt{\_Decimal64}, and \texttt{\_Decimal128} extension types. All of these are C
libraries: they are compiled and linked, expose no \texttt{constexpr} interface, and
do not run on accelerators.

On the C++ side, ISO/IEC TR 24733 \cite{tr24733} and the later WG21 proposals N3407
and N3871 \cite{wg21-n3407, wg21-n3871} specified decimal class types for C++, but
proposed to implement them by wrapping the existing Intel or IBM C libraries and
treated \texttt{constexpr} support as allowed rather than required; no proposal has
been adopted into the standard. Portable header-only numerical libraries in the
broader Boost and TOMS tradition, such as the e\_float multiple-precision system
\cite{algorithm910}, MPFR \cite{mpfr}, and the customized floating-point framework
FloatX \cite{floatx}, demonstrate the value of a header-only, standards-integrated
C++ approach, but none provides IEEE 754 decimal floating-point.

\subsection{Decimal on GPUs}
There is, to our knowledge, no prior IEEE 754 decimal floating-point implementation
that runs on a graphics processor. The closest existing work is the fixed-point
decimal support in NVIDIA's RAPIDS cuDF library, whose \texttt{decimal32},
\texttt{decimal64}, and \texttt{decimal128} column types are fixed-point (a runtime
scale, in the manner of Apache Spark and Java BigDecimal), not floating-point, and
whose arithmetic manipulates scales rather than a self-adjusting exponent
\cite{nvidia-cudf}. The IEEE 754 support in NVIDIA GPU hardware is binary only.
Decimal floating-point on the GPU is therefore new ground, which we describe in
Section~\ref{sec:gpu}.

\subsection{Algorithmic foundations}
Boost.Decimal does not invent new arithmetic; it assembles well-established algorithms in a
single portable, \texttt{constexpr} C++ form, and cites each where it is used. The
wide-integer multiply and divide underneath the significand arithmetic are the schoolbook
Algorithm M and Algorithm D of Knuth \cite{knuth_taocp2}, with the inner quotient step using
the reciprocal method of Moller and Granlund \cite{moller_granlund2011}. Division and
rounding by powers of ten use reciprocal multiplication by an invariant divisor
\cite{granlund_montgomery1994, warren_hd}. Integer-to-text conversion uses the JEAIII method
\cite{jeaiii}, and shortest round-trip output uses Ryu \cite{ryu}. The square root follows
the table-and-Newton approach of Berkeley SoftFloat \cite{softfloat}, adapted to base ten.
The elementary functions combine Cody-Waite range reduction \cite{cody_waite1980}, minimax
polynomial approximations \cite{muller2016}, the arithmetic-geometric mean for elliptic
integrals \cite{dlmf, algorithm910}, an accelerated series for the Riemann zeta function
\cite{borwein1995}, and Kahan compensated summation \cite{kahan1965}.
Table~\ref{tab:algorithms} collects these. The contribution of this work is not the
algorithms individually but their delivery together as one portable, header-only type.

\begin{table}
  \caption{Algorithmic foundations: the established methods the implementation uses, by
    component. The novelty of this work is their unified, portable, \texttt{constexpr}
    delivery, not the methods themselves.}
  \label{tab:algorithms}
  \footnotesize
  \begin{tabular}{lll}
    \toprule
    Component & Method & Reference \\
    \midrule
    Wide multiplication       & Schoolbook (Algorithm M)         & \cite{knuth_taocp2} \\
    Wide division             & Algorithm D; reciprocal 3/2      & \cite{knuth_taocp2, moller_granlund2011} \\
    Divide or round by $10^k$ & Reciprocal by invariant divisor  & \cite{granlund_montgomery1994, warren_hd} \\
    Integer to text           & JEAIII                           & \cite{jeaiii} \\
    Shortest float to text    & Ryu                              & \cite{ryu} \\
    Square root               & Table plus Newton (SoftFloat)    & \cite{softfloat} \\
    Range reduction           & Cody-Waite                       & \cite{cody_waite1980} \\
    Polynomial approximation  & Minimax (Remez)                  & \cite{muller2016} \\
    Elliptic integrals        & Arithmetic-geometric mean        & \cite{dlmf, algorithm910} \\
    Riemann zeta              & Accelerated alternating series   & \cite{borwein1995} \\
    Series summation          & Kahan compensated sum            & \cite{kahan1965} \\
    \bottomrule
  \end{tabular}
\end{table}

\subsection{Summary of the gap}
Table~\ref{tab:comparison} summarizes how Boost.Decimal differs from the established
implementations. Each individual property has precedent; the contribution of this
work is to provide all of them at once.

\begin{table}
  \caption{Boost.Decimal compared with the established decimal floating-point
    software. Each property has precedent individually; providing all of them
    together is the contribution of this work.}
  \label{tab:comparison}
  \footnotesize
  \begin{tabular}{lllll}
    \toprule
     & \textbf{Boost.Decimal} & \textbf{Intel BID} & \textbf{IBM decNumber} & \textbf{GCC \texttt{\_Decimal}} \\
    \midrule
    Language           & C++14                 & C          & C           & C extension \\
    Header-only        & yes                   & no         & no          & no \\
    \texttt{constexpr} & yes                   & no         & no          & no \\
    GPU / CUDA         & yes                   & no         & no          & no \\
    Architectures      & x86/ARM/s390x/PPC64LE & x86 only   & portable C  & GCC targets \\
    License            & BSL-1.0               & BSD-style  & ICU         & GPL + RLE \\
    \bottomrule
  \end{tabular}
\end{table}

%% ===================================================================
\section{Design and Architecture}
\label{sec:design}

\subsection{Goals}
The library is guided by three goals, in priority order: ease of use, portability, and
performance. Ease of use means the types behave like the built-in floating-point types
and integrate with the standard library, so that adopting decimal is largely a matter
of changing a type name. Portability means a single source compiles and produces
identical results on every supported target without per-platform code paths.
Performance means decimal arithmetic should be competitive with, and where possible
faster than, the established software implementations. This section describes the type
set, the internal representation, and the \texttt{constexpr} strategy that ties the
first two goals together; portability and performance are treated in
Sections~\ref{sec:portability} and \ref{sec:performance}.

\subsection{The decimal type set}
Boost.Decimal provides two families of types. The conforming family,
\texttt{decimal32\_t}, \texttt{decimal64\_t}, and \texttt{decimal128\_t}, implements
the IEEE 754 interchange formats exactly, including their bit-level encoding, so that a
value can be stored, transmitted, and recovered across implementations. The second
family, \texttt{decimal\_fast32\_t}, \texttt{decimal\_fast64\_t}, and
\texttt{decimal\_fast128\_t}, computes the same numerical results but trades storage
width for speed. Both families present the same interface, following the operator and
function set proposed in ISO/IEC TR 24733 \cite{tr24733}. Mixed-type arithmetic follows
rules in the spirit of the usual arithmetic conversions: an operation between two
decimal types yields the higher-precision type, and an operation between a conforming
type and its fast counterpart yields the fast type. Table~\ref{tab:formats} lists the
three conforming formats and their parameters.

\begin{table}
  \caption{The conforming decimal formats. Each packs a sign, an exponent, and an integer
    significand into a word of the given width using the BID encoding.}
  \label{tab:formats}
  \footnotesize
  \begin{tabular}{lrrl}
    \toprule
    Type & Storage & Digits $p$ & Significand range \\
    \midrule
    \texttt{decimal32\_t}  & 32-bit  & 7  & $0$ to $10^{7}-1$ \\
    \texttt{decimal64\_t}  & 64-bit  & 16 & $0$ to $10^{16}-1$ \\
    \texttt{decimal128\_t} & 128-bit & 34 & $0$ to $10^{34}-1$ \\
    \bottomrule
  \end{tabular}
\end{table}

\subsection{Internal representation}
A conforming type stores a single unsigned integer holding the IEEE 754 BID encoding;
\texttt{decimal32\_t}, for instance, is a class wrapping one \texttt{std::uint32\_t}.
The sign, a combination field, the biased exponent, and the integer significand are all
packed into those bits, so each operation must decode its operands, compute, and
re-encode the result. A fast type instead stores the significand, exponent, and sign in
three separate, already-normalized fields: \texttt{decimal\_fast32\_t}, for example,
holds an unsigned significand, a small unsigned exponent, and a boolean sign. This
removes both the per-operation decode and the normalization that the conforming types
need in order to compare across cohorts (the several encodings of one value, such as 1e5
and 10e4, that decimal permits). The cost is that fast types occupy more memory, carry
no cohort information, and flush subnormals to zero; the benefit is a measurable speedup,
quantified in Section~\ref{sec:performance}. The fast types are not, however, uniformly
faster: producing the shortest output, for instance, requires stripping the trailing
zeros that normalization introduces.

\subsection{The constexpr strategy}
A central decision is that the entire library is usable in a constant expression.
C++14's relaxed \texttt{constexpr} rules make this possible: with loops and local
mutation permitted in \texttt{constexpr} functions, arithmetic, parsing, formatting, and
the elementary functions can all be evaluated at compile time, which lets the types
stand in for built-in types in template metaprogramming and in contexts that require
constant initialization. This choice recurs throughout the implementation. It precludes
reliance on the C floating-point environment for exception flags (discussed below), and
it shapes the code so that the same function bodies are valid both at compile time and
on a GPU device, where the CUDA compiler eagerly type-checks and emits device code for
every instantiation. The library detects whether an evaluation is occurring at compile
time and selects the matching rounding-mode source, so that constant evaluation and run
time agree.

\begin{lstlisting}[caption={Decimal arithmetic in a constant expression: the sum is formed
  and checked entirely at compile time.},label={lst:constexpr}]
using boost::decimal::decimal64_t;
constexpr decimal64_t tenth {1, -1};                          // 0.1
static_assert(tenth + tenth + tenth == decimal64_t{3, -1});   // 0.3, exactly
\end{lstlisting}

\subsection{Standard-library integration and deviations}
The types are accompanied by decimal versions of the standard headers a numeric type is
expected to support, including \texttt{<cmath>}, \texttt{<charconv>},
\texttt{<cstdlib>}, \texttt{<cfenv>}, \texttt{<limits>}, \texttt{<format>}, and the
user-defined literals. Meeting the C++14 baseline and the conceptual requirements of
Boost.Math additionally makes the types usable with that library's statistics,
quadrature, and interpolation facilities without extra work. A few deliberate deviations
follow from the design goals and are documented as such. Because the library chooses
\texttt{constexpr} over the C floating-point environment, it does not set or read the
floating-point exception flags, which would otherwise require the C-style out-parameter
signatures used by the Intel library. The elementary functions are not claimed to be
correctly rounded to half a unit in the last place, matching both the behavior and the
stated position of the reference C libraries. Input and output add a
\texttt{cohort\_preserving\_scientific} format so that a value's cohort can be printed,
\texttt{from\_chars} reports overflow and underflow distinctly (following the resolution
direction of LWG issue 3081) and honors the active rounding mode, and parsing of
non-finite values from a stream is permitted. We return to the I/O deviations in
Section~\ref{sec:io}.

%% ===================================================================
\section{Arithmetic Implementation}
\label{sec:impl}

Every finite decimal value is a signed integer significand scaled by a power of ten.
The four basic operations decode their operands into this significand, exponent, and
sign form, compute an exact intermediate result in an integer wide enough to hold it,
and then pass through a shared finalization step that shrinks the intermediate to the
type's precision $p$ (7, 16, or 34 digits for the three sizes), rounds it, and packs the
value back into the result type. Finalization is also where a result that is too large
becomes an infinity and one that is too small flushes to zero. For results known to be
in range the finalizer writes the BID bit pattern directly and skips the general
constructor's bounds checks; out-of-range results take the constructor path.

\subsection{Addition and subtraction}
Addition first aligns the operands to a common exponent by scaling the significand with
the smaller exponent up by the appropriate power of ten, then adds or subtracts the
aligned significands according to their signs. The result can exceed the target
precision by a few digits, or, after a large alignment shift, by many. The common case
of one to three excess digits is detected by a chained comparison against precomputed
thresholds, while the rare larger case is handled by counting digits and looking up the
divisor in a power-of-ten table. Either way the excess is removed by a single division
with round half to even; division by a constant power of ten is performed by reciprocal
multiplication by an invariant divisor~\cite{granlund_montgomery1994, warren_hd} rather
than a hardware divide. Subtraction is addition with the sign of the second operand
inverted.

\subsection{Multiplication}
Multiplication adds the operand exponents and multiplies the significands by the
schoolbook method, Knuth's Algorithm M~\cite{knuth_taocp2}. With both
significands expanded to full precision $p$, the exact product has $2p-1$ or $2p$
digits and is held in an integer twice the significand width: a 64-bit product for the
32-bit types, a 128-bit product for the 64-bit types, and a 256-bit product for the
128-bit types. The finalizer branches on the threshold $10^{2p-1}$ to decide whether to
drop $p-1$ or $p$ digits, divides by the corresponding power of ten with round half to
even, handles the carry that arises when rounding lifts the quotient to $10^{p}$, and
packs the result.

\subsection{Division}
Division subtracts the operand exponents and forms the quotient by long division of the
significands (Knuth's Algorithm D~\cite{knuth_taocp2}), with the inner quotient-digit
estimate using the reciprocal 3/2 method of Moller and Granlund~\cite{moller_granlund2011}.
The dividend significand is pre-scaled by $10^{p}$ so that dividing by the
divisor significand yields a quotient of exactly $p$ or $p+1$ digits together with a
remainder, and that remainder is precisely the sticky information needed to round
correctly. With both operands at full precision the quotient lies in
$[10^{p-1}, 10^{p+1})$; the finalizer applies round half to even using the remainder,
and the dropped low digit when the quotient has $p+1$ digits, again handling the carry
to $10^{p}$. The 32-bit case keeps the scaled dividend in a 64-bit integer, with the
wider backend types used for the larger precisions.

\subsection{Rounding and special values}
The default rounding direction is round to nearest, ties to even, which the inline
finalizers above assume; the other IEEE 754 rounding directions are honored by routing
through a separate step that reads the active mode. Because the library is
\texttt{constexpr}, it distinguishes a compile-time evaluation, where the rounding mode
is a translation-time constant, from a run-time evaluation, where the mode is read
through the \texttt{<cfenv>} interface, so that the two agree. Signed zeros, infinities,
and quiet and signaling NaNs follow the IEEE 754 rules; for example, subtracting two
like-signed infinities and multiplying zero by an infinity each produce a quiet NaN.
Comparison of the conforming types normalizes cohorts so that numerically equal values
compare equal regardless of encoding, whereas the fast types, being stored normalized,
compare directly.

\subsection{The wide-integer backend}
The arithmetic above relies on integers wider than 64 bits. Rather than depend on a
compiler's \texttt{\_\_int128}, which is absent on 32-bit targets and on MSVC, the
library bundles a self-contained backend providing 128-bit signed and unsigned types and
a 256-bit type, with the multiplication and division primitives they require; on targets
that offer a native 128-bit type the backend uses it. This single backend serves the
wide products of multiplication, the scaled dividends of division, and the square-root
computation below, and it is a principal reason one source compiles identically on every
supported architecture (Section~\ref{sec:portability}).

\subsection{Square root}
\label{sec:sqrt}
Square root illustrates how an elementary function is implemented entirely in integer
arithmetic, so that it is usable in a constant expression and produces identical results
on every target, given that no hardware decimal square root exists. The method follows
Berkeley SoftFloat \cite{softfloat}, adapted from radix two to radix ten. The argument is
reduced to a value $g \in [1, 10)$ while its exponent is tracked separately; an initial
approximation of $1/\sqrt{g}$ is read from a 90-entry table that covers $[1, 10)$ in
steps of $0.1$, with linear interpolation between adjacent entries; the approximation is
refined by Newton-Raphson iterations expressed as integer remainder corrections of the
form $z' = z + (x - z^2)/(2z)$; and the result is rescaled by the appropriate power of
ten, with a further factor of $\sqrt{10}$ when the original exponent is odd. The number
of refinement iterations and the integer width grow with precision, as
Table~\ref{tab:sqrt} shows: the 128-bit case uses the 256-bit backend for the squaring
and the remainder and rounds to nearest, while the smaller types round by floor.
Recasting the computation in this integer form made it several times faster than the
library's earlier square root, by roughly a factor of six to eight for the 32- and
64-bit types and about four for the 128-bit type.

\begin{table}
  \caption{Square root by precision: the integer width of the working arithmetic, the
    number of Newton-Raphson refinement iterations, and the final rounding rule.}
  \label{tab:sqrt}
  \footnotesize
  \begin{tabular}{llll}
    \toprule
    Type & Integer width & Newton iterations & Final rounding \\
    \midrule
    \texttt{decimal32\_t}  & 64-bit  & 1 & floor \\
    \texttt{decimal64\_t}  & 128-bit & 2 & floor \\
    \texttt{decimal128\_t} & 256-bit & 3 & round to nearest \\
    \bottomrule
  \end{tabular}
\end{table}

\subsection{Special functions}
\label{sec:special}
Beyond the basic operations, the library supplies a \texttt{constexpr} \texttt{<cmath>}
surface that goes well past the standard's mandate, including the exponential, logarithmic,
and power functions, the trigonometric and hyperbolic functions, the error and gamma
family, and several orthogonal-polynomial and elliptic-integral families. The
implementations use textbook methods adapted to fixed decimal precision: Cody-Waite range
reduction \cite{cody_waite1980} to bring an argument into a small interval, minimax (Remez)
polynomial approximations \cite{muller2016} on that interval, the arithmetic-geometric mean
for the complete elliptic integrals \cite{dlmf, algorithm910}, an accelerated alternating
series for the Riemann zeta function \cite{borwein1995}, and Kahan compensated summation
\cite{kahan1965} where a long series would otherwise lose accuracy. They build on the
multiple-precision special-function techniques of Algorithm 910 \cite{algorithm910}, by one
of the present authors. Like the reference libraries they are not correctly rounded
(Section~\ref{sec:design}); we do not count them among the contributions of this paper, but
they show that the types form a complete numeric facility rather than a minimal arithmetic
core.

%% ===================================================================
\section{Exact Decimal I/O and Conformant Conversion}
\label{sec:io}

For a decimal type the boundary between text and value is not an afterthought; it is the
main reason to use the type. A price such as 19.99, entered as text, must become the
value nineteen and ninety-nine hundredths exactly, and must print back as 19.99, with no
binary rounding introduced on the way in or out. This section describes how
Boost.Decimal parses and prints values, how it exposes the IEEE 754 encodings, and where
its conversion interface deviates, by design, from the C++ standard.

\subsection{Shortest round-trip output}
The default text output produces the shortest decimal string that reads back to the same
value, using the Ryu algorithm of Adams \cite{ryu} generalized to the 128-bit
significand of \texttt{decimal128\_t}. Because the stored value is already base ten, the
shortest-representation problem is simpler than it is for binary types, and the printed
digits are exact rather than the nearest binary approximation. The conversion of the
significand digits to text uses the JEAIII integer-to-string method \cite{jeaiii}. Parsing
and printing are offered through a \texttt{<charconv>}-style interface,
\texttt{from\_chars} and
\texttt{to\_chars}, which returns a pointer and a \texttt{std::errc} status and, like the
rest of the library, is usable in a constant expression. A companion compile-time trait
reports the largest buffer each format and precision can require, so callers can size
output buffers exactly instead of guessing.

\subsection{Formats and cohorts}
The interface provides the scientific, fixed, hexadecimal, and general formats of the
standard, plus one addition specific to decimal: a cohort-preserving scientific format. A
cohort is the set of distinct encodings of one value that decimal admits; in a
seven-digit type the value three hundred can be written as any of 3e+2 through
3000000e-4, each a different bit pattern denoting the same number. By default the
conversion functions ignore cohorts and emit the normalized form, which is the behavior
users expect from a numeric type. When a caller needs to record exactly which member of
the cohort a value is, the cohort-preserving format reads and writes the significand and
exponent verbatim, so that the round trip is bitwise and not merely numerically faithful.
The format carries the constraints the standard implies: it is unavailable for the fast
types, which store no cohort, and combining it with an explicit output precision is
rejected, because rounding or zero-padding would destroy the very cohort information the
format exists to preserve.

\subsection{The IEEE 754 bit encodings}
Decimal values can also be exchanged at the bit level. The library exposes conversions to
and from both IEEE 754 significand encodings: the Binary Integer Decimal (BID) form it
uses internally, and the Densely Packed Decimal (DPD) form used by hardware decimal
units. The \texttt{to\_bid}, \texttt{from\_bid}, \texttt{to\_dpd}, and \texttt{from\_dpd}
functions move a value to or from a fixed-width unsigned integer of the matching size, so
a value produced by Boost.Decimal can be handed to hardware or to another implementation
and recovered exactly. These conversions are likewise \texttt{constexpr} and run on a GPU
device.

\subsection{Conformance and deliberate deviations}
The \texttt{from\_chars} and \texttt{to\_chars} interfaces follow the standard
\texttt{<charconv>} functions where possible, and the library provides overloads taking
\texttt{std::chars\_format} that return the standard result types, so existing code ports
with little change. Three deviations are intentional and follow from the nature of
decimal. First, \texttt{from\_chars} distinguishes overflow from underflow, by setting
the output value, where the standard reports only a single out-of-range error; this
follows the direction of LWG issue 3081 and matches Boost.Charconv. Second, because the
decimal types are far more sensitive to rounding direction than binary types are,
\texttt{from\_chars} rounds according to the active rounding mode rather than always to
nearest. Third, the cohort-preserving format above has no counterpart in the standard.
Each deviation is documented, and the standard-compatible overloads remain available for
code that requires strict conformance.

%% ===================================================================
\section{Portability and Cross-Platform Consistency}
\label{sec:portability}

Portability in Boost.Decimal means more than compiling on many systems: it means that
the same program produces the same bits on every one of them. This section describes how
that is achieved and how it is verified.

\subsection{One source, no per-platform arithmetic}
The arithmetic is written once, in portable C++, and does not branch on the target
architecture. Two mechanisms make this possible. The first is the bundled wide-integer
backend of Section~\ref{sec:impl}: where a native 128-bit integer exists the library uses
it, and where it does not, on 32-bit targets and on MSVC, the emulated type takes over
with identical behavior, so multiplication and division need no architecture-specific
code. The second is explicit control of byte order. Because a conforming value is a
packed bit pattern, its bytes must be interpreted consistently regardless of the host's
endianness; the library determines endianness at compile time and refuses to build, with
a request to file an issue, on any target whose byte order it cannot establish, rather
than silently producing wrong results. This is what lets one BID encoding be correct on
little-endian x86 and on big-endian s390x alike.

\subsection{Tolerating compiler and standard-version differences}
A header-only library that supports compilers as old as GCC 8 and Clang 6 must absorb a
wide range of compiler behavior without changing its results. Boost.Decimal does so by
detecting language and library features and degrading gracefully: it uses
\texttt{if constexpr} where available and a plain conditional under C++14, a
\texttt{constexpr} bit-cast where the compiler supports one and a fallback where that
support is known to be broken (for instance on a particular Clang release and on 32-bit
targets), and a working detection of constant evaluation across the several compiler
versions that implement it differently. The C++14 floor is itself a portability choice:
it is the lowest standard on which the \texttt{constexpr} design can be expressed, and it
admits the largest set of toolchains.

\subsection{The tested matrix}
The library is built and tested across the axes listed in Table~\ref{tab:platforms},
spanning 32- and 64-bit word sizes, little- and big-endian byte orders, and hosted and
bare-metal targets, with the major compiler on each. Continuous integration builds and
runs the full test suite, including the conformance suite of
Section~\ref{sec:verification}, on this matrix on every change, and the tests assert
identical numerical results across it, so a regression on any one architecture,
endianness, or word size is caught at once. It is this cross-platform agreement, verified
rather than assumed, that makes the library suitable where reproducibility across
machines is a requirement.

\begin{table}
  \caption{Portability coverage. The library is built and tested across these axes;
    PPC64LE and the bare-metal target are exercised under QEMU emulation.}
  \label{tab:platforms}
  \footnotesize
  \begin{tabular}{ll}
    \toprule
    Dimension & Coverage \\
    \midrule
    Operating systems & Linux, macOS, Windows \\
    Architectures     & x86-32, x86-64, ARM32, ARM64, s390x, PPC64LE \\
    Word size         & 32-bit and 64-bit \\
    Byte order        & little-endian and big-endian (s390x) \\
    Bare metal        & ARM Cortex-M (no operating system) \\
    Compilers         & GCC 8+, Clang 6+, Visual Studio 2019+, Intel oneAPI \\
    \bottomrule
  \end{tabular}
\end{table}

%% ===================================================================
\section{GPU and Accelerator Support}
\label{sec:gpu}

A distinguishing capability of Boost.Decimal is that the same decimal types run on an
NVIDIA GPU. Because the established libraries are host C code, decimal arithmetic has not
previously been available inside CUDA kernels. This section describes how the library
provides it and what is tested.

\subsection{One source for host and device}
The \texttt{constexpr} discipline of Section~\ref{sec:design} is what makes device
support practical. Every function the library exposes is annotated with a macro that
expands to nothing for an ordinary host build and to the CUDA
\texttt{\_\_host\_\_ \_\_device\_\_} qualifiers when the translation unit is compiled by
the CUDA compiler with device support enabled. Because the function bodies were already
written to be valid in a constant expression, they contain no constructs, no exceptions
and no run-time-only library calls, that would prevent the device compiler from accepting
and emitting them, so a single source serves both host and device with no separate device
implementation to maintain.

\subsection{Device execution}
On the device the library makes a few automatic adjustments that the host build does not
require: it uses the bundled emulated 128-bit integer rather than a compiler built-in,
since device support for native 128-bit integers is limited, and it disables exceptions
and the assertion macro, which have no place in device code. The numerically meaningful
behavior is unchanged, so a decimal computation performed in a kernel produces the same
result as the same computation on the host. The arithmetic, the comparisons, the BID and
DPD conversions, the mathematical constants, and text conversion through
\texttt{<charconv>} are all available on the device, across the three precisions and in
mixed-precision combinations.

\subsection{Testing on the device}
Device support is validated by a dedicated CUDA test suite. For each of
\texttt{decimal32\_t}, \texttt{decimal64\_t}, and \texttt{decimal128\_t}, and for mixed
combinations of them, kernels exercise construction, the four arithmetic operations,
comparison, the BID and DPD encodings, the constants, and \texttt{<charconv>}; inputs are
transferred and results read back through CUDA managed memory, and each device result is
checked against the host computation. Passing this suite is the evidence behind the claim
that the library computes identically on the GPU and on the CPU. Table~\ref{tab:gpu}
summarizes what the device suite covers.

\begin{table}
  \caption{Operations validated on the CUDA device, by precision. Each device result is
    checked against the host result. Mixed-precision combinations of the three types are
    also exercised.}
  \label{tab:gpu}
  \footnotesize
  \begin{tabular}{lccc}
    \toprule
    Operation & \texttt{decimal32\_t} & \texttt{decimal64\_t} & \texttt{decimal128\_t} \\
    \midrule
    Construction              & yes & yes & yes \\
    Add, subtract, multiply, divide & yes & yes & yes \\
    Comparison                & yes & yes & yes \\
    BID encode and decode     & yes & yes & yes \\
    DPD encode and decode     & yes & yes & yes \\
    Mathematical constants    & yes & yes & yes \\
    Text conversion (charconv) & yes & yes & yes \\
    \bottomrule
  \end{tabular}
\end{table}

\subsection{Scope of the novelty claim}
We are careful about the scope of this claim. To our knowledge no prior IEEE 754 decimal
floating-point library runs on a GPU. The nearest existing capability, the decimal column
types in NVIDIA's cuDF, is fixed-point rather than floating-point
(Section~\ref{sec:background}), and the hardware's own IEEE 754 support is binary. We
therefore state the contribution as the first decimal floating-point library with GPU
support that we are aware of, rather than as an absolute first, and we note that the
library does not at present provide a non-CUDA accelerator path such as SYCL.

%% ===================================================================
\section{Conformance and Verification}
\label{sec:verification}

Correctness is the first requirement of a numerical library, and for a portable one it
must hold identically on every target. Boost.Decimal is checked at four levels: against
the official decimal arithmetic conformance suite, by property-based and fuzz testing of
the most input-exposed code, by a body of regression tests, and by the cross-platform
agreement of Section~\ref{sec:portability}. Together these form the evidence on which a
Replicated Computational Results review (Appendix~\ref{app:rcr}) would rest.

\subsection{Conformance to IEEE 754}
The conforming types implement the IEEE 754 decimal formats, their arithmetic, the
rounding directions, and the special values. Two deviations are documented and
deliberate, both discussed in Section~\ref{sec:design}: the library does not maintain the
floating-point exception flags, having chosen compile-time evaluation over the C
floating-point environment, and the elementary functions are not claimed to be correctly
rounded to half a unit in the last place, which matches the behavior and the stated
position of the reference C libraries. The mandated operations, including the basic
arithmetic and square root, are otherwise implemented to the standard's requirements.

\subsection{The decimal conformance suite}
The primary correctness check is the General Decimal Arithmetic test suite, the
machine-readable conformance vectors that accompany Cowlishaw's specification
\cite{Cowlishaw2003} and against which the reference implementations are themselves
tested. The suite is vendored into the repository and driven by a parser and runner that
reads each vector, performs the named operation in the appropriate type, and compares the
result and status against the expected values. The vectors cover the core operation set,
including absolute value, addition, subtraction, multiplication, division, the comparison
operations and their signaling and total-order variants, maximum and minimum, quantize,
and remainder, in both the 64-bit and 128-bit formats. Table~\ref{tab:dectest} groups the
covered operations. Running the same external suite that the Intel and IBM libraries use is
what lets us claim conformance rather than merely internal consistency.

\begin{table}
  \caption{Operations covered by the vendored General Decimal Arithmetic conformance
    vectors. The vectors are supplied as the 64-bit (dd) and 128-bit (dq) test sets, with a
    subset for the 32-bit format, in 35 files in total.}
  \label{tab:dectest}
  \footnotesize
  \begin{tabular}{ll}
    \toprule
    Category & Operations \\
    \midrule
    Arithmetic & add, subtract, multiply, divide \\
    Comparison & compare, compare-signaling, compare-total-order \\
    Rounding   & quantize, remainder \\
    Other      & absolute value, negate, maximum, minimum, clamp \\
    \bottomrule
  \end{tabular}
\end{table}

\subsection{Property-based and fuzz testing}
The text-conversion surface is the most exposed to untrusted input and the most prone to
edge cases, so it is fuzzed continuously. A set of libFuzzer harnesses drives parsing and
formatting through every entry point, including \texttt{from\_chars},
\texttt{strtod}-style parsing, the string constructors, and the \texttt{printf},
\texttt{<format>}, and \{fmt\}-style output paths, each seeded with a corpus and run in
continuous integration. Round-trip properties, that printing a value and reading it back
recovers the value, and that the cohort-preserving format recovers the exact bit pattern,
are checked directly.

\subsection{Regression and cross-platform checks}
Every bug fixed during development is captured as a regression test so that the same
defect cannot reappear; there are at present roughly forty such tests, alongside the
operation- and function-level unit tests. Crucially, the entire suite, the conformance
vectors, the property tests, and the regression tests, is run on every target in the
portability matrix of Table~\ref{tab:platforms}, and the results are required to agree bit
for bit. Line coverage is measured and reported in continuous integration. A defect that
appeared only on a big-endian or a 32-bit target, or only during constant evaluation,
would be caught by this matrix rather than discovered in the field.

%% ===================================================================
\section{Performance Evaluation}
\label{sec:performance}

This section quantifies how Boost.Decimal compares with the established software
implementations of decimal arithmetic. The headline question is not whether decimal is
slower than binary, it is, because binary floating-point has hardware support on every
target of interest while all decimal arithmetic here is software, but whether
Boost.Decimal is competitive with the reference decimal software: the Intel library and
the GCC built-in types.

\subsection{Methodology}
Each arithmetic benchmark fills a vector of twenty million random values and applies the
operation between successive elements, repeating the measurement five times to obtain a
stable figure; the comparison benchmark applies the six relational operators, and the
text benchmarks perform twenty million conversions. The same harness measures the
built-in \texttt{double}, the Boost.Decimal conforming and fast types, the GCC
\texttt{\_Decimal} built-ins, and the Intel BID library. One caveat applies throughout:
the Intel and GCC types are driven through C routines that are close to, but not identical
with, the C++ routines used for the Boost.Decimal types, because those types are C
constructs; we take the routines to be near enough for a fair comparison, as their
authors' own benchmarks do. The complete per-operation and per-platform tables are in the
library documentation; we report a representative subset here.

\subsection{Platforms}
Results were collected on four configurations: x86-64 Linux on an Intel i9-11900k with
both the Intel oneAPI compiler and GCC 13.4; x86-64 Windows on the same processor with
Visual Studio; ARM64 macOS on an Apple M4 Max with Clang; and ARM64 Windows. The trends
below hold across all four; we draw the table from the x86-64 Linux GCC configuration
because it measures every baseline at once.

\subsection{Results}
Table~\ref{tab:perf} reports the time to perform twenty million 64-bit operations on
x86-64 Linux. Against the Intel reference library, Boost.Decimal is faster on every
operation, by factors ranging from roughly two for multiplication to about ten for
comparison, and the fast type widens the margin further. Against the GCC
\texttt{\_Decimal64} built-in, the fast type is faster for addition, subtraction, and
division, and comparable for multiplication and comparison. The conforming
\texttt{decimal64\_t}, which additionally pays for cohort handling, is competitive with
the GCC type and remains well ahead of the Intel library.

\begin{table}
  \caption{Time in microseconds to perform twenty million 64-bit operations on x86-64
    Linux (Intel i9-11900k, GCC 13.4); lower is better. All four columns are software
    implementations of decimal arithmetic. For reference, hardware binary
    \texttt{double} completes the same arithmetic workloads in well under one hundred
    thousand microseconds.}
  \label{tab:perf}
  \footnotesize
  \begin{tabular}{lrrrr}
    \toprule
    Operation & \texttt{decimal64\_t} & \texttt{decimal\_fast64\_t} & GCC \texttt{\_Decimal64} & Intel \texttt{BID\_UINT64} \\
    \midrule
    Comparison     & 1,090,026 & 513,587   & 475,605   & 5,452,325 \\
    Addition       & 1,257,515 & 1,064,025 & 2,100,094 & 3,379,104 \\
    Subtraction    & 1,313,266 & 1,166,909 & 1,313,261 & 3,522,244 \\
    Multiplication & 2,528,577 & 2,312,590 & 2,462,813 & 4,973,224 \\
    Division       & 2,477,436 & 1,724,535 & 2,904,760 & 5,745,255 \\
    \bottomrule
  \end{tabular}
\end{table}

The advantage is largest at 128 bits, where the reference library is weakest: on x86-64
Linux with the Intel compiler, multiplying \texttt{decimal\_fast128\_t} is about
3.7 times faster than the Intel \texttt{BID\_UINT128}, and on Windows the multiplication
gap is wider still. With the Intel compiler the fast 64-bit addition is about 4.6 times
faster than \texttt{BID\_UINT64}, somewhat better than the GCC-toolchain figure in the
table, which reflects the compiler difference rather than a change in the library.
Table~\ref{tab:perf128} gives the 128-bit operations under the GCC toolchain. It also shows
the one case where a competitor leads: at 128-bit multiplication the GCC built-in is
modestly ahead of Boost.Decimal, though both remain far ahead of the Intel library.

\begin{table}
  \caption{Time in microseconds to perform twenty million 128-bit operations on x86-64
    Linux (Intel i9-11900k, GCC 13.4); lower is better. All four columns are software
    implementations of decimal arithmetic.}
  \label{tab:perf128}
  \footnotesize
  \begin{tabular}{lrrrr}
    \toprule
    Operation & \texttt{decimal128\_t} & \texttt{decimal\_fast128\_t} & GCC \texttt{\_Decimal128} & Intel \texttt{BID\_UINT128} \\
    \midrule
    Comparison     & 3,772,309 & 444,145   & 872,026   & 5,705,342 \\
    Addition       & 2,085,250 & 926,772   & 3,264,420 & 2,509,923 \\
    Multiplication & 7,966,781 & 8,116,908 & 6,683,052 & 21,352,000 \\
    Division       & 5,854,061 & 4,442,123 & 9,906,049 & 14,918,326 \\
    \bottomrule
  \end{tabular}
\end{table}

For text conversion the picture changes, because binary \texttt{<charconv>} is also
software: here the decimal types are competitive with or faster than \texttt{double}. On
x86-64 Linux, formatting at a fixed six-digit precision, \texttt{decimal32\_t} and
\texttt{decimal64\_t} cost roughly half as much as \texttt{double} in scientific notation
and well under it in general notation, and parsing is within a small factor of
\texttt{double} across the formats.

\subsection{Summary}
Across the basic operations and on every platform measured, Boost.Decimal is faster than
the Intel reference library, and it is faster than the GCC built-in decimal types in many
operations though not all; for text conversion it is competitive with or faster than
hardware-backed binary \texttt{double}. We state the result deliberately as parity or
improvement over the reference software, not as a universal speed record: against hardware
binary floating-point, decimal arithmetic remains several times slower, the expected cost
of exact base-ten behavior.

%% ===================================================================
\section{Applications: Composability with Boost.Math}
\label{sec:applications}

The applications payoff of Boost.Decimal is not that it reproduces base-ten arithmetic;
that follows from the design and was shown in Sections~\ref{sec:io} and
\ref{sec:portability}. It is that, because the types are standards-conforming C++ types
rather than an opaque C library, they compose with the existing C++ numerical ecosystem. A
large body of analysis becomes available on exact decimal data with no new code, and with
the same cross-platform reproducibility as the arithmetic. This is something the Intel and
IBM C libraries cannot offer, because they are not C++ types and do not model the concepts
that generic numerical code is written against.

\subsection{Exact input}
We use a year of daily stock prices as a running dataset. Read with the decimal types, the
data is exact: summing 252 closing prices in \texttt{decimal32\_t} reproduces the
spreadsheet total of 52151.99 to the cent, whereas the same sum in \texttt{float}
accumulates to 52151.96. The error is not in the addition but in the representation, and it
is removed at the point of parsing. With the values exact, any subsequent analysis starts
from the right numbers.

\subsection{Statistics for free}
Because the decimal types satisfy the user-defined-type requirements of Boost.Math, that
library's algorithms operate on a \texttt{std::vector} of decimals directly. The same calls
one would write for \texttt{double}, the mean, median, and variance functions of
\texttt{boost::math::statistics}, compute the closing-price statistics, and a square root
from the library's own \texttt{<cmath>} gives the standard deviation, from which two-sigma
Bollinger bands follow. On the sample year this yields a mean closing price of \$207.20, a
median of \$214.26, a standard deviation of \$25.45, and bands at \$258.11 and \$156.30. No
decimal-specific statistics code was written: the analysis is the generic Boost.Math
algorithm instantiated on the decimal type. In practice a user enables one documented macro
so that the mixed integer-decimal expressions inside the generic code compile, and
suppresses a few of Boost.Math's float-comparison warnings; the algorithms themselves are
unchanged.

\begin{lstlisting}[caption={Boost.Math's statistics algorithms run on a vector of decimals
  with no decimal-specific code; the Bollinger band follows from the mean and the standard
  deviation.},label={lst:stats}]
using boost::decimal::decimal64_t;
namespace stats = boost::math::statistics;

std::vector<decimal64_t> closes = /* parsed exactly from the CSV */;
const decimal64_t mean = stats::mean(closes);
const decimal64_t var  = stats::variance(closes);
const decimal64_t sd   = boost::decimal::sqrt(var);
const decimal64_t upper = mean + 2 * sd;   // upper two-sigma Bollinger band
\end{lstlisting}

\subsection{The breadth, and why it comes for free}
The statistics above are one corner of what the same concept conformance unlocks. The
library's continuous integration exercises the decimal types with Boost.Math across a much
wider surface: the univariate statistics suite in full (the first four moments, variance
and sample variance, skewness, kurtosis, median and median absolute deviation, the
interquartile range, the Gini coefficient, and the mode), and a range of special functions
(the beta and Riemann zeta functions, the elliptic integrals, and the Legendre, Laguerre,
and Hermite polynomial families), with the library's own elementary functions cross-checked
against Boost.Math's. The rest of the Boost.Math toolkit, its statistical distributions,
quadrature, interpolation, and root-finding, is reachable through the same mechanism, since
all of it is generic code written against the same numeric concepts.

The reason this comes essentially for free is generic programming. A Boost.Math algorithm
is written once against a concept and instantiated on any type that models it; meeting the
C++14 baseline and those concept requirements (Section~\ref{sec:design}) is therefore
sufficient to make the toolkit available on decimal values, and the results inherit the
bit-for-bit cross-platform consistency of Section~\ref{sec:portability}. A C library of
decimal routines, however complete, cannot be dropped into this generic code, because it
provides functions rather than a type that models the concepts. The composability is thus
not an add-on but a direct consequence of implementing decimal as a first-class C++ type,
and it is, for a numerical-software audience, among the most consequential practical
results of the design.

%% ===================================================================
\section{Discussion and Limitations}
\label{sec:discussion}

The design choices behind Boost.Decimal carry costs as well as benefits, and an honest
account of them sharpens the contribution.

\subsection{The conforming and fast types}
The fast types buy speed with space and standards conformance. They are wider than the
conforming types, they do not represent cohorts (so they cannot preserve a value's exact
encoding through I/O), and they flush subnormal values to zero. They are also not
uniformly faster: producing the shortest decimal string, for instance, requires removing
the trailing zeros that their normalized storage introduces, which can make text output
slower than for the conforming types. A user therefore chooses per use: the conforming
types where the IEEE 754 bit layout, cohorts, or subnormals matter, and the fast types
where throughput dominates.

\subsection{The cost relative to binary}
Decimal arithmetic in Boost.Decimal is software, and on every platform of interest binary
floating-point is hardware; the basic operations are accordingly several times to a few
tens of times slower than \texttt{double} (Section~\ref{sec:performance}). This is inherent
to decimal rather than to this implementation, and it bounds the library's applicability:
it is the right tool where exact base-ten behavior or cross-platform reproducibility is
required, and the wrong one for bulk numerical work that is indifferent to base. The
elementary functions are likewise not correctly rounded to half a unit in the last place,
matching the reference C libraries but falling short of a correctly-rounded \texttt{libm}.

\subsection{Accelerator and platform scope}
GPU support is at present specific to CUDA; there is no SYCL or other portable accelerator
backend, and on the device the library uses its emulated 128-bit integer because native
support is limited. The tested platform matrix is broad but not exhaustive, and the
cross-platform consistency guarantee is only as strong as that matrix.

\subsection{Threats to validity in the comparison}
The performance comparison has the limitations its methodology implies. The Intel and GCC
baselines are driven through C routines that are close to, but not identical with, the C++
routines used for the decimal types, so small differences may reflect the harness rather
than the arithmetic. Results were collected on a few machines rather than averaged over
many, and they depend on the compiler and its optimization settings, as the gap between
the Intel-compiler and GCC figures in Section~\ref{sec:performance} shows. We therefore
present the comparison as evidence of parity or better with the reference software, not as
a precise speedup factor.

%% ===================================================================
\section{Conclusions and Future Work}
\label{sec:conclusion}

Boost.Decimal shows that a single, portable, header-only C++ source can provide IEEE 754
decimal floating-point that runs at compile time, on the CPU, on a GPU, and on bare-metal
targets, with results identical across architectures and performance that matches or
exceeds the established reference software. The formats and the encoding are not new; the
contribution is to deliver them together, as a first-class C++ type, without the
integration cost, the loss of \texttt{constexpr}, or the platform limits of the existing C
libraries. Because the type models the concepts that generic numerical code expects, the
broader C++ ecosystem, Boost.Math among it, operates on exact decimal data unchanged, which
turns a single arithmetic type into a foundation for reproducible decimal analysis.

Several directions remain. A correctly-rounded mode for the elementary functions would close
the one conformance gap the library shares with the reference C libraries. A portable
accelerator backend, such as SYCL, would extend device support beyond CUDA. The tested
platform matrix can be broadened further, and the Boost.Math composability of
Section~\ref{sec:applications} invites larger case studies in the domains, finance and
commerce above all, where exact decimal analysis matters most. The library is released under
the Boost Software License and developed in the open.

%%
\begin{acks}
  We thank The C++ Alliance for sponsoring the development of this work, and the Boost
  community reviewers and the review manager, John Maddock, for the peer review that
  preceded the library's acceptance into Boost.
\end{acks}

%%
%% Bibliography. Uses the ACM-Reference-Format style from the template folder
%% and the references.bib alongside this file.
\bibliographystyle{ACM-Reference-Format}
\bibliography{references}

%%
\appendix
\section{Reproducibility (for the RCR review)}
\label{app:rcr}

All results in this paper can be reproduced from the public repository. This appendix lists
what to build and how, so a reviewer can replicate the conformance results and the
benchmarks.

\paragraph{Source.}
The results were produced from a tagged commit of the repository at
\url{https://github.com/boostorg/decimal}; the exact tag is recorded with the artifact. The
library is header-only and requires only a C++14 compiler.

\paragraph{Tests and conformance.}
The full test suite, including the General Decimal Arithmetic conformance vectors of
Section~\ref{sec:verification}, is built and run with Boost.Build from the library's
\texttt{test} directory:
\begin{verbatim}
b2 cxxstd=20 toolset=gcc-13
\end{verbatim}
Running the same command under other toolsets (for example \texttt{toolset=clang} or
\texttt{toolset=msvc}) and on the architectures of Table~\ref{tab:platforms} reproduces the
cross-platform agreement; targets without native hardware are exercised under QEMU.

\paragraph{Benchmarks.}
The benchmarks are compiled into the test programs and enabled with a preprocessor
definition. From the \texttt{test} directory:
\begin{verbatim}
b2 cxxstd=20 toolset=gcc-13 \
   define=BOOST_DECIMAL_RUN_BENCHMARKS benchmarks -a release
\end{verbatim}
Appending \texttt{,BOOST\_DECIMAL\_BENCHMARK\_CHARCONV=1} to the definition also runs the
text-conversion benchmarks. The reference baselines build directly from their C sources:
the GCC built-in decimal baseline with
\begin{verbatim}
gcc benchmark_libdfp.c -O3 -std=c17
\end{verbatim}
and the Intel BID baseline by linking the Intel library,
\begin{verbatim}
icx benchmark_libbid.c -O3 $PATH_TO_LIBBID/libbid.a -std=c17
\end{verbatim}
each producing an executable that prints per-operation timings.

\paragraph{Hardware and reading the results.}
The figures in Section~\ref{sec:performance} were collected on an Intel i9-11900k for the
x86-64 results and an Apple M4 Max for the ARM64 macOS results. Each benchmark reports the
total time in microseconds for twenty million operations, repeated five times for
stability, so lower numbers are better, and the ratios between types are the quantities of
interest rather than the absolute times.

\end{document}
